\cotefile{100}
\scanpage{2}
Bonne \unsure{réduction} et gpes $p$-divisibles
% archive sheet number: 1

\scanpage{3}
\mspage{1}
\uline{Théorème} $S$ schéma loc. noeth, $U$ ouvert, $T$ fermé complémentaire, $p$ \unsure{nb|est} premier, tel que $\forall t\in T$, $\mathrm{car}\ k(t)=p$. Considérons la catégorie $C$ des schémas abéliens \inserted{$A$} sur $S$, et la catégorie $C'$ des couples $(A_U,M,\varphi)$, où $A_U$ est un schéma abélien sur $U$, $M$ un \unsure{groupe|gpe} $p$-divisible sur $S$, et $\varphi:M|U\longrightarrow T_p(A_U)$. Soit
\[
F:C\longrightarrow C'
\]
le foncteur
\[
F(A)=(A\times_S U,T_p(A),\mathrm{id}_{T_p(A)|U}).
\]

Alors :

1) Supposons que $U$ rencontre toute composante connexe de $S$ \inserted{(p.ex. $U$ dense dans $S$, i.e. $\operatorname{prof\,ét}_T(S)\geq 1$).} Alors le foncteur $F$ est pleinement fidèle.

2) Supposons que \deleted{$U$ soit dense dans $S$ et que} a) \deleted{\unsure{le morphisme} \illegible{2} \unsure{de} $S$ \unsure{sur} $S$, \unsure{et}} \inserted{$\exists$ morphisme fini surjectif. $S'\longrightarrow S$ avec $S'$ géom. unibranche (p.ex. la normalisée de $S$ soit finie sur $S$, \unsure{où|ou|si} $S$ local).}

\deleted{b) pour tout $t\in T$, si $S'$ est le localisé strict de $S$ en $t$, \inserted{$\bar t$ son point fermé,} et $S'_1$ et $S'_2$ deux composantes irréductibles de $S'$, on \deleted{\unsure{ait}} ait $S'_1\cap S'_2\neq\{\bar t\}$. \deleted{Sous ces conditions, $F$ est une équivalence de catégories.} (N.B. Cette condition équivaut aussi à celle-ci : l'image inverse $U''$ de $U$ dans $S''=S'\times_S S'$ est dense dans $S''$. Elle est vérifiée si $S$ est \deleted{\unsure{géom.}} unibranche en les points de $T$) Alors \unsure{a)} b) Posant} b) Que \unsure{-|.} $\operatorname{prof\,ét}_T(S)\geq 2$\tikz[baseline=0pt,x=1ex,y=1ex]{\draw (0,2.1)--(.06,.20) .. controls (-.03,-.12) and (-.43,-.09) .. (-.28,.16) .. controls (-.2,.31) and (-.04,.27) .. (.06,.20);}. Alors $F$ est \deleted{pleinement fidèle} \inserted{une équivalence de} catégories

\margin{Pour 1) TSVP : $\longrightarrow$}

\tikz[baseline=0pt,x=1ex,y=1ex]{\draw (0,2.1)--(.06,.20) .. controls (-.03,-.12) and (-.43,-.09) .. (-.28,.16) .. controls (-.2,.31) and (-.04,.27) .. (.06,.20);} Soit, avec les notations \deleted{de} \inserted{du} N° a), $S''=S'\times_S S'$, $S'''=S'\times_S S'\times_S S'$. Considérons l'ensemble simplicial $\overset{\substack{L_0\\\text{''}}}{\pi_0(S')}\leftleftarrows\overset{\substack{L_1\\\text{''}}}{\pi_0(S'')}\mathrel{\substack{\leftarrow\\[-.65ex]\leftarrow\\[-.65ex]\leftarrow}}\overset{\substack{L_2\\\text{''}}}{\pi_0(S''')}$, \deleted{\unsure{les|Com}} et \deleted{\unsure{soit}} soit $K_0\leftleftarrows K_1\mathrel{\substack{\leftarrow\\[-.65ex]\leftarrow\\[-.65ex]\leftarrow}} K_2$ le ss-ensemble simplicial formé avec les composantes connexes dont l'image dans $S$ rencontre $U$. Supposons \unsure{:} que : \deleted{pour tout} \inserted{b')} la plus petite partie de $L_1$ contenant $K_1$ et stable par ``troisième face'' soit $L_1$ (c'est ce qui est assuré par la condition b)). Alors $F$ est \deleted{\unsure{pl. fid.|pl. f.}} une équivalence de catégories .....

% archive sheet number: 2

\scanpage{4}
\mspage{1 bis}
\uline{Lemme} $A,B$ schémas abéliens sur $S$ loc. noeth, $u_\ell:T_\ell(A)\unsure{-|\longrightarrow}T_\ell(B)$. Alors $\exists$ \uline{sous-schéma fermé} $T$ de $S$ tel que, pour tout $S'/S$, on \inserted{\unsure{ceci|ait ceci|ait} :} $u_{\ell S'}$ provient de $u:A_{S'}\unsure{-|\longrightarrow}B_{S'}$ \unsure{sss|ssi} $S'$ se factorise par $T$.

Cf Inventions 1.2 (\unsure{ou|où|une} démonstration \unsure{cristalline|conditionnelle}).

\deleted{\uline{Cor} Supposons que $\forall s\in\unsure{S|\operatorname{Ass}\mathcal O_S}$, \unsure{l'une des|l'image de} hyp. suiv}

\uline{Complément} De plus, \deleted{\unsure{de la conclusion}} pour tt $t\in T$ tel que $\mathrm{car.}\,k(t)=\ell$, $T$ majore un voisinage ouvert immédiat de $t$.

\uline{Corollaire} Supposons que l'ouvert $U\subset S$ soit tel que

a) tte comp. conn. de $S$ rencontre $U$

b) $t\in S-U\Longrightarrow\mathrm{car.}\,k(t)=\ell$.

c) $u_\ell|U$ provient de $u_U:A/U\unsure{-|\longrightarrow}B/U$

Alors \deleted{\unsure{toute}} $u_\ell$ provient de $u:A\longrightarrow B$

\deleted{En effet, on peut supposer $S$ connexe, \inserted{\unsure{donc} a) $\Longleftrightarrow U\neq\varnothing$.} Alors \deleted{$\unsure{T\text{ est }S}$} \unsure{tout pt de} \deleted{\unsure{l'une des} composantes irréductibles} [distinguer $t\in\deleted{\unsure{T}}U$ et $t\in T\cap(S-U)$ \inserted{grâce à b)}], \deleted{\unsure{de} $S$. Si on avait \unsure{disons} $T\neq S$ on aurait une comp. $S_i$ de $S$ telle} Comme $T\neq\varnothing$, il suffit de prouver que tte comp. irréd. $S_i$ de $S$ qui rencontre $T$ est $\subset T$. Or si $S_i\cap U=\varnothing$ \unsure{clair|clairement}, \inserted{car $T\supset\overline U$} \unsure{sinon} aussi grâce au complément \inserted{et à b)}. Donc $T\supset S_{\mathrm{red}}$. \deleted{\unsure{Ainsi}} Mais $T$ est \deleted{\unsure{indu|fermé}} ouvert} $T$ est ouvert \unsure{évident|induit} sur $U$ par hyp. \inserted{\deleted{et} c)}, et en les pts de $(S-U)\cap T$ grâce au complément et à b). Donc $\deleted{U}\unsure{T|\overline U|\overline T}$ est un \inserted{sous-schéma} ouvert, et comme il est fermé, c'est une réunion de comp. connexes. Par a) on a donc $T=S$, cqfd

% archive sheet number: 3

\scanpage{5}
\mspage{2)}
\deleted{\unsure{2}}

\[
T\unsure{\subset|\leftarrow}S\supset U=S-T
\]
\hspace*{1em}car. résiduelles toutes $=p>0$.

$S$ noethérien, $U$ \deleted{dense} \deleted{\unsure{schémati|schématiquement}} \deleted{schématiquement} dense

\deleted{\unsure{2)}}
\diagram{\begin{tikzcd}[ampersand replacement=\&,column sep=large,row sep=1ex]
 A\arrow[r,mapsto] \& (A_U,T_p^{\deleted{\unsure{p|\natural|\varphi}}}(A),\varphi)\\
 \substack{\text{schémas}\\\text{abélien}\\\text{sur }S} \& \substack{\text{triples }(A_U,M,\varphi)\\ A_U\ \text{schéma abélien sur }U\\ M\ \text{gpes }p\text{-div sur }T\\ \varphi:T_p^{\deleted{p}}(A_U)\simeq M.}
\end{tikzcd}}

\margin{Pour pleine fidélité \deleted{\unsure{si p ex. $S$ irréd. ou car resp $\neq\ell$}} \deleted{\unsure{Au}} lieu de $U$ dense, \inserted{et $S-U$ de car. $p$} il \deleted{\unsure{ne}} suffit \unsure{(?)|que} que \deleted{\unsure{pour}} tte comp. \deleted{\unsure{irréd.}} \deleted{\unsure{de $S$} \illegible{3} $\unsure{T}$ \unsure{non vide} $U$} \deleted{et $\overline U$ est de complément \unsure{de} $S$} \deleted{\unsure{de chaque}} connexe de $S$ rencontre $U$, \deleted{p.ex. $S$} \deleted{pleine fidélité} (\deleted{pour l'effectivité} sinon contre exemple avec $U$ \inserted{connexe et} non vide \deleted{\unsure{+}} obtenu en attachant deux spectres d'anneaux de val. discrète complets, et $U=$ un des deux pts maximaux \ldots) Y-a-t-il des contre'exemples avec $S$ de \unsure{étages|Noether|dim.} \unsure{non équi|non égaux|$2$} ?}

\uline{1)} \uline{Foncteur pleinement fidèle}

fidèle bien connu, car $A\unsure{-|\longrightarrow}T_{\deleted{\unsure{\ell}}p}^{\deleted{p}}(A)$ est fidèle.

``pl.'' fidèle : il faut prouver que si $u_p:T_p^{\deleted{p}}(A)-T_p^{\deleted{p}}(B)$ est algébrique sur $U$, il l'est sur $S$. Si $U$ est sch. dense dans $S$, cela résulte de $[\,]$ 1.2. Dans le cas général, soit $X_0$ l'adhérence schématique de $U$ dans $X$. Alors on dispose d'un hom $\varphi_0:A_{X_0}\longrightarrow B_{X_0}$, \deleted{\unsure{On veut|Comme}} \inserted{compatible avec $u_p$.}

\deleted{\deleted{Soit $\mathcal J$ l'Idéal qui définit $X_0$ sur $X$, il est nilpotent et \unsure{concentré} sur $T$.} Alors les \unsure{théories|thèses|th.} de Serre-Tate impliquent que l'on peut prolonger la restriction de $u_0$ à $T_{\deleted{\mathrm{red}}0}$ (\inserted{$=T$} muni de la structure induite réduite) \deleted{en} un hom. $A_{T_n}\longrightarrow B_{T_n}$, où $T_n$ est le $n$\textsuperscript{ième} voisinage infinitésimal. Se ramenant au cas $S$ affine, \inserted{$=\operatorname{Spec}\Lambda$} on trouve \deleted{que l'on} un hom $A_{\widehat\Lambda}-B_{\widehat\Lambda}$, $\widehat\Lambda=$ complété de $\Lambda$ pour la top. définie par \deleted{\unsure{le}} l'idéal de $T_0$ dans $S_0$. Comme \unsure{:} l'hom $U\times_{\operatorname{Spec}\widehat\Lambda}\longrightarrow\operatorname{Spec}\Lambda$ est fid plat, on gagne .....} car appliquer alors $[\,]$ 1.1.

% archive sheet number: 4

\scanpage{6}
\mspage{-0.13-}
image of $\handwritten{\phi'}$. Since this image is dense, there are no such $y$'s. Therefore $\handwritten{\sigma}^{\handwritten{\unsure{\prime|\vee}}}[\handwritten{\phi}^{-1}(\handwritten{\mathcal U})]=\handwritten{\phi}^{-1}(\handwritten{\mathcal U})\times_{\handwritten{\mathcal U}}\handwritten{\mathcal U'}$, and $\handwritten{\phi'}$ is surjective. Therefore $\handwritten{\overline\omega}$ is geometrically injective without restriction, and by \uline{Zariski's Main Theorem}, $\handwritten{\overline\omega}$ is an isomorphism. \uline{QED}
% Pencil underline beneath Zariski's Main Theorem.
\margin{\handwritten{!}}

\uline{Remark} \uline{(5)}. The following are equivalent:

(a) $(Y,\handwritten{\phi})$ is a universal categorical quotient of $X$ by $G$,

(b) for all affine schemes $Y\handwritten{'}$, and morphisms $Y\handwritten{'\longrightarrow}Y$, if $\handwritten{\phi'}:X\handwritten{'\longrightarrow}Y\handwritten{'}$ is the base extension, then $(Y\handwritten{'},\handwritten{\phi'})$ is a categorical quotient of $X\handwritten{'}$ by $G$,

(c) there is an affine open covering $\handwritten{\{}\handwritten{\mathcal U}_i\handwritten{\}}$ of $Y$ such that if $\handwritten{\phi}_i:\handwritten{\phi}^{-1}(\handwritten{\mathcal U}_i)\handwritten{\longrightarrow}\handwritten{\mathcal U}_i$ is the restriction of $\handwritten{\phi}$, then $(\handwritten{\mathcal U}_i,\handwritten{\phi}_i)$ is a universal categorical quotient of $\handwritten{\phi}^{-1}(\handwritten{\mathcal U}_i)$ by $G$, for all $i$.

(6). In Chapter 1, we will need a criterion for $(Y,\handwritten{\phi})$ to be a categorical quotient which does not imply that $(Y,\handwritten{\phi})$ is a geometric quotient. In fact, the following 3 conditions suffice:

i) $\handwritten{\phi}\circ\handwritten{\sigma}=\handwritten{\phi}\circ p_2$,

ii) $\underline{\mathrm{o}}_Y$ is the subsheaf of invariants of $\handwritten{\phi}_*(\underline{\mathrm{o}}_X)$,

iii) if $W$ is an invariant closed subset of $X$, then $\handwritten{\phi}(W)$ is closed in $Y$; if $W_i$, $i\in I$, is a set of invariant closed subsets of $X$, then:
\[
\handwritten{\phi}\bigl(\underset{i\in I}{\cap}\,W_i\bigr)=\underset{i\in I}{\cap}\,\handwritten{\phi}(W_i)
\]
% slanted
\margin{\handwritten{\diagram{\begin{tikzcd}[ampersand replacement=\&,row sep=small,column sep=small]
 R\arrow[d,shift right,"p_1"']\arrow[d,shift left,"p_2"] \&\\
 \varphi:X\arrow[r,"f"'] \& Y
\end{tikzcd}} morphism of preschemes \unsure{tel que|s.t.} $fp_1=fp_2$ \unsure{makes|induces|we say} $Y$ \unsure{into|is} a quotient $X/R$ if (ii) $\mathcal O_Y$ is the sheaf of invariants of $\mathcal O_Y-f_*(\mathcal O_X)\rightrightarrows g_*(\mathcal O_{\unsure{X|R}})$ (iii) Any invariant open $U$ of $X$ is the inverse image of an open subset of $Y$}}
\margin{\handwritten{\unsure{It point suffices if|I point suffices if|I think sufficient if} $\varphi$ \uline{is qu. cpct} qu. sep.}}

% archive sheet number: 5

\scanpage{7}
\mspage{3}
2) \uline{Foncteur ess. surjectif} \deleted{\unsure{est}} \inserted{\uline{si} $S$ \uline{géom. unibranche en les points de} $S-U$ \unsure{(ou|(} \inserted{[plus gén, si, posant $S''=S'\times_S S'$, $U'=U\times_S S'$, $U''=U\times_S S''$, $U'\deleted{\unsure{\text{est dense}}}$ \inserted{$(U'')$} rencontre tte composante connexe de $S'$ $(S'')$]} et si \unsure{le normalisé|la normalisée|la normalisation} $S'$ de $S_{\mathrm{red}}$ est \unsure{fini|finie} sur $S$.}

\deleted{a) Supposons}

Partons de $A_U,M,\varphi$, on veut construire $A$. On peut supposer $S$ affine, $S=\operatorname{Spec}\Lambda$. \deleted{Pour \unsure{descendre|avoir|descente} $A$, il suffit de traiter le cas}

a) Réduction au cas $S$ réduit.

Soit. $S_0=S_{\mathrm{réd}}$, $A_0$ schéma abélien sur $S_0$ qui y résout le pb. Par Serre-Tate, on trouve \deleted{des} \deleted{\unsure{un}} schémas abéliens sur les $T_n$, donc un schéma \inserted{formel} abélien \inserted{$A'$} sur \deleted{\unsure{\ell}} $S'=\operatorname{Spec}\widehat\Lambda$. Son image inverse sur $S'_0$ ``est'' un vrai schéma abélien, \deleted{sur $\unsure{S_0}$} savoir l'image inverse de $A_0$ sur $S'_0$.
\diagram{\begin{tikzcd}[ampersand replacement=\&,row sep=small,column sep=small]
 S_0\arrow[d] \& S'_0\arrow[l]\arrow[d]\\
 S \& S'\arrow[l]
\end{tikzcd}}
Donc (ref \ldots\ !!!) $A'$ provient d'un vrai schéma abélien $A'$ sur $S'$ [l'imm. $S'_0-S'$ étant \unsure{nilpotente}]. Or $U\amalg S'\longrightarrow S$ est fid. plat, et on trouve une donnée de descente sur $A_U\amalg A'$. Elle devient effective après chgt de base $S_0\unsure{\longleftarrow|-|\longrightarrow}S$, donc elle l'était déjà avant, d'après la \unsure{section|critère|notion|réduction} de Raynaud. On gagne !

\margin{?}
\deleted{b) Réduction au cas $S$ local, \deleted{\unsure{et pt}} $T$ son point fermé. Pour faire cette réduction, on est ramené à prouver ceci : si $t$ est un pt \unsure{fermé\ldots?} de $T$, $V$ un ouvert \unsure{cont|contenant} $t$ et $U$, \deleted{et} $A_V$ un schéma abélien sur $V$ \deleted{qui résout} tel que sa restriction à $S'=\operatorname{Spec}\mathcal O_{S,t}$ \deleted{résolve} soit OK., alors $A_V$ est \deleted{\unsure{encore}} OK dans un voisinage, i.e. l'isom donné $T_p(A_{S'})\simeq M_{S'}$ se prolonge en un isom en un}

% archive sheet number: 6

\scanpage{8}
\mspage{-0.14-}
\uline{Proof}: Note first of all that $\handwritten{\phi}$ is dominating by (ii); hence by (iii), $\handwritten{\phi}$ is actually \uline{surjective}. Now suppose $\deleted{\handwritten{\phi}}\inserted{\handwritten{\Psi}}:X\handwritten{\longrightarrow} Z$ is any $S$-morphism such that $\handwritten{\psi}\circ\handwritten{\sigma}=\handwritten{\psi}\circ p_2$. We proceed exactly as in the proof of Proposition 0.1: choose an affine open covering $\handwritten{\{}V_i\handwritten{\}}$ of $Z$. Once we show that there is an open covering $\handwritten{\{}\handwritten{\mathcal U}_i\handwritten{\}}$ of $Y$ such that \tikz[baseline=(v.base)]{\node[draw,ellipse,inner xsep=2pt,inner ysep=1pt] (v) {$\handwritten{\psi}^{-1}(V_i)\handwritten{\supset}\handwritten{\phi}^{-1}(\handwritten{\mathcal U}_i)$};}, we can conclude the proof as in Proposition 0.1. To show this, let $W_i=X-\handwritten{\psi}^{-1}(V_i)$. These are closed invariants subsets of $X$. Therefore, by (iii), $\handwritten{\mathcal U}_i=Y-\handwritten{\phi} W_i$ are open; and $\handwritten{\phi}^{-1}(\handwritten{\mathcal U}_i)\handwritten{\subset}\handwritten{\psi}^{-1}(V_i)$ for all $i$. Finally, $\handwritten{\{}\handwritten{\psi}^{-1}(V_i)\handwritten{\}}$ cover $X$, hence $\underset{i}{\cap}\,W_i=\handwritten{\phi}$, hence by (iii) $\underset{i}{\cap}\,\handwritten{\phi} W_i=\handwritten{\phi}$, hence $\handwritten{\{}\handwritten{\mathcal U}_i\handwritten{\}}$ cover $Y$. \uline{QED}

(7). Suppose $(Y,\handwritten{\phi})$ is a geometric quotient of $X$ by $G$. Notice that conditions (i), (ii), and (iii) of Definition 0.6 are preserved by any base extension $Y^{\handwritten{\prime}}\handwritten{\longrightarrow} Y$. Only the more delicate condition (iv) might (and indeed sometimes does) break down. However, if $Y^{\handwritten{\prime}}\handwritten{\longrightarrow} Y$ is \uline{flat}, then condition (iv) is preserved. To see this, let $\handwritten{\psi}=\handwritten{\phi}\circ\handwritten{\sigma}=\handwritten{\phi}\circ p_2$. Then (iv) asserts the exactness of:
% slanted
\margin{\handwritten{$\varphi$ quasi-compact quasi-séparated}}
\[
0\handwritten{\longrightarrow} \underline{\mathrm{o}}_Y\handwritten{\longrightarrow} \handwritten{\phi}_*(\underline{\mathrm{o}}_X)\xrightarrow{\handwritten{\sigma}^*-p_2^*}\handwritten{\psi}_*(\underline{\mathrm{o}}_G\times_S X)
\tag*{$(*)$}
\]
But taking direct image sheaves commutes with \uline{flat} base extension, i.e.
\[
\begin{aligned}
\handwritten{\phi}_*(\underline{\mathrm{o}}_X)\handwritten{\otimes}\underline{\mathrm{o}}_Y &\mathrel{\overset{\handwritten{\sim}}{=}}\handwritten{\phi}^{\handwritten{\prime}}_*(\underline{\mathrm{o}}_X)\\
\handwritten{\psi}_*(\underline{\mathrm{o}}_{G\times_S X})\handwritten{\otimes}\underline{\mathrm{o}}_Y &=\handwritten{\psi}^{\handwritten{\prime}}_*(\underline{\mathrm{o}}_{G\times_S X^{\handwritten{\prime}}}).
\end{aligned}
\]
(where $X^{\handwritten{\prime}}=X\times_Y Y^{\handwritten{\prime}}$, etc.), and the exactness of a sequence is pre-

% archive sheet number: 7

\scanpage{9}
\deleted{voisinage ouvert $W\ni t$, induisant ce qui est donné sur $W\cap U$. En fait, \deleted{d'après le} \inserted{si on dispose du} th. de Tate \deleted{\unsure{ceci}} \unsure{pas seulmt|pas seult} en \unsure{inégales|inégale} caractéristique, on pourrait prouver}

b) Cas $S$ \deleted{normal} \inserted{géom. unib.}, $A_U$ muni d'une rigidification de Jacobi. On a donc un morphisme
\margin{\unsure{R}}
\[
U\longrightarrow R
\]
\margin{schéma modulaire à échelon [utilisant aussi une polarisation de $A_U$]}
d'où par \inserted{l'adhérence du} le graphe un sous-schéma fermé $S'$ de $S\times R$. Le critère valuatif et le résultat de Raynaud dans le cas d'un trait impliquent que $S'\longrightarrow S$ est propre. Il reste à prouver que c'est un isom [tout revient en effet, grâce à Tate, à prouver que $A_U$ se prolonge en $A$, quand $T_p(A_U)$ se prolonge en un $M$ -- cet $M$ étant en effet unique à iso unique près). Par le Main Theorem, on est ramené à prouver que \deleted{les pts de $S'$} \inserted{ttes fibres} $S'_s$ \unsure{sont} \uline{discrètes}. Introduisons la normalisée $S''$ de $S'$, qu'on supposera d'abord \uline{finie} sur $S'$. Alors le Th. de Tate implique que le gpe $\unsure{b|p}$-divisible $T_p(\underbrace{A_{S''}}_{\inserted{\text{image inverse de }A_R}})$ est l'image inverse de $\deleted{\unsure{T_p\text{ de }A|T_p(A)}}M$. Donc la restriction de $A_{S''}|S''_s$, correspondant à un $T_p$ ``constant'', et $S''_s$ étant géom connexe,
\diagram{\begin{tikzcd}[ampersand replacement=\&,row sep=small,column sep=small]
 \& S''\arrow[d] \& U''\arrow[l,hook]\arrow[d,phantom,"\rotatebox{90}{$\simeq$}" description]\\
 R \& S'\arrow[l]\arrow[d] \& U'\arrow[l,hook]\arrow[d,phantom,"\rotatebox{90}{$\simeq$}" description]\\
 \& S \& U\arrow[l,hook]
\end{tikzcd}}

% archive sheet number: 8

\scanpage{10}
\mspage{-0.15-}
served under flat base extension -- hence the analogous sequence $(*)^{\handwritten{\prime}}$ on $Y^{\handwritten{\prime}}$ is exact.

(8). On the other hand, consider the converse. I claim that if $Y^{\handwritten{\prime}}\handwritten{\longrightarrow} Y$ is faithfully flat and quasi-compact, and if the extended morphism $\handwritten{\phi}^{\handwritten{\prime}}:X^{\handwritten{\prime}}\handwritten{\longrightarrow} Y^{\handwritten{\prime}}$ makes $(Y^{\handwritten{\prime}},\handwritten{\phi}^{\handwritten{\prime}})$ a universal categorical quotient or geometric quotient or universal geometric quotient of $X^{\handwritten{\prime}}$ by $G$, then $(Y,\handwritten{\phi})$ is the same type of quotient of $X$ by $G$. These assertions are all simple corollaries of the general theory of descent of SGA 8.

\fbox{\S3.} We shall now look somewhat more closely at the structure of an action of a group. As above, let $G/S$ act via $\handwritten{\sigma}$ on $X/S$.

\uline{Definition} \uline{0.8}: The action $\handwritten{\sigma}$ is said to be

i) \uline{closed} if for all geometric points $x$ of $X$, the orbit $O(x)\handwritten{\subset}\overline X$ is closed, (i.e. $\overline X=X\times_S\operatorname{Spec}\handwritten{\Omega}$ if $x$ is an $\handwritten{\Omega}$-valued point),

ii) \uline{separated} if the image of
\[
\handwritten{\unsure{\Psi|\overline{\underline\psi}}}=(\handwritten{\sigma},p_2):G\times_S X\handwritten{\longrightarrow} X\times_S X
\]
is closed,

iii) \uline{proper} if $\handwritten{\Psi}$ is proper,

iv) \uline{free} if $\handwritten{\Psi}$ is \tikz[baseline=(v.base)]{\node[draw,ellipse,inner xsep=2pt,inner ysep=1pt] (v) {a closed immersion.};}
% slanted
\margin{\handwritten{\fbox{monomorphisme}}}
% slanted
\margin{\handwritten{conflits avec des terminologies reçues dire : \uline{proper and free}}}

Moreover, let $\handwritten{\mathscr S}=S(1_X)$ be the stabilizer of $1_X:X\handwritten{\longrightarrow} X$ (cf. Definition 0.4). It is a sub-group pre-scheme of the group pre-scheme $G\times_S X/X$. Let $\handwritten{\overline\omega}:\handwritten{\mathscr S}\handwritten{\longrightarrow} X$ be the projection and let $e_X:X\handwritten{\longrightarrow}\handwritten{\mathscr S}$

% archive sheet number: 9

\scanpage{11}
on conclut (cf lemme) que $A_{S''}$ est constant. \deleted{Il l'est aussi en tant que variété polarisée, par fonctorialité. Il ne reste le choix que d'un nb fini de rigidifications de Jacobi, \unsure{donc} $(S''_s)_{\mathrm{réd}}\longrightarrow R$ se factorise par $\unsure{\mathrm{Sp}| \operatorname{Spec}}(k(s))$, donc $S''_s$ est réduit à un point.} On gagne.

Le cas où on ne suppose plus $S''$ fini sur $S'$ se traite par \unsure{la méthode|les méthodes} essentiellement de $[\,]$ p 75-76, \uline{à condition} qu'on prouve que l'isomorphisme sur $U''\simeq U$ entre $T_p(A'')|U''$ et $M_{S''}|U''$ se prolonge à $S''$. \deleted{Il faudrait tirer au clair cette question.} \inserted{OK}
\margin{?}

c) Cas $S$ \deleted{normal} \inserted{géom. unibranche}. Choisissons un \unsure{ent|ent.} premier $\ell\neq p,3$, et soit $U'\longrightarrow U$ un morphisme étale fini surjectif \deleted{qui \unsure{rig}} tel que $A_{U'}$ soit muni d'une \unsure{rigigidification|rigidification} d'échelon $\ell$.

% slanted
\margin{\unsure{utilisons|initialisons|inutile} un \unsure{rev.|fond} !}
Normalisons $S$ dans $R(U')/R(U)$ \inserted{$R(S)$}, d'où morphisme fini \inserted{surjectif} $S'\longrightarrow S$, \inserted{prolongeant $U'-U$} avec $S'$ normal intègre. Alors on peut trouver un schéma abélien $A'$ sur $S'$ qui prolonge $(A_U)\times_U U'$, avec $T_p(A')\simeq p^*(M)$.

\diagram{\begin{tikzcd}[ampersand replacement=\&,row sep=small,column sep=small]
 U\arrow[d,arrows={Hooks[right]-}] \& U'\arrow[l]\arrow[d,arrows={Hooks[right]-}] \& U''\arrow[l,shift left]\arrow[l,shift right]\arrow[d,arrows={Hooks[right]-}]\\
 S \& S'\arrow[l,"p"'] \& S''\arrow[l,shift left]\arrow[l,shift right]
\end{tikzcd}}
Alors par la partie ``pl. fid'' appliquée à $S''=S'\times_S S'$, utilisant que $U''$ est dense dans $S''$ (\deleted{\unsure{par Chevalley}} \inserted{N.B.} $S''-S$ est ouvert car $S'-S$ est univ\textsuperscript{t} ouvert \ldots), on trouve

% archive sheet number: 10

\scanpage{12}
\mspage{2}
Errata, Page 2.

$A/\handwritten{\unsure{\mathfrak p|\mathscr P}}^{(k)}$ in the expression for $E^m$. Now $A/\handwritten{\mathfrak p}^{(k)}$ has depth 1, whether calculated as an $A$-module or a $B$-module. So it has homological dimension $r-1$ over $B$, since $B$ is regular, and hence each $\operatorname{Ext}^r_B(A/\handwritten{\mathfrak p}^{(k)},B)=0$. Thus the $E^r$ term of the abutment is zero, so each ''

Page 43, Line 4: insert ``locally free $F$, and hence for all'' after ``for all''.

Page 46, Line 9: insert ``proper over $k$, '' before ``and''.

Page 46, Line -2: add ``[AVB 2.2 and 3.3]'' at end of line.

Page 50, Line -6: replace ``locally a complete intersection'' by ``non-singular''.

Page 51, Lines -6, -5: ditto

Page 64, Lines 5, 6, 7 read instead: `` (b) if $\mathrm{tr.d.}\ K(A)/k=d+1$, then $K(A)$ is a finitely generated extension of $K$''.

Page 70, lines -2, -1 and Page 71, Line 1: ditto

Page 73, Lines 11, 12, 13: put a period after ``$K(\hat X)$,'' and delete the following three lines, beginning at ``i.e. ''.

Page 75, Lines 2, 3, 4: replace this sentence by ``The only points where $G$ does not act freely are the two lines $X_0=X_2=0$ and $X_1=X_3=0$, hence $Z$ is non-singular except at the images of these two lines.

Page 75, Line 10: delete $y_2^2,\ y_2 y_3^2,\ y_3^4$.

Page 75, Line 11: change 12 to 4.

% archive sheet number: 11

\scanpage{13}
que $A'$ est muni d'une \uline{donnée de descente} relativement à $S'\unsure{-|\to}S$. En particulier, \deleted{\unsure{l|E}} $ {}_{\ell}A'$ est muni d'une telle donnée de descente,. \deleted{qui} \deleted{est} Or $S'-S$ est \inserted{morphisme} de descente effective pour les schémas relatifs finis étales (SGA 1 VIII), donc on trouve un \deleted{\unsure{se|tel}} tel \unsure{rev.|sch.|schéma} \inserted{en groupes} \deleted{\unsure{B|R}} $N$ de $S$, tel que $N\mid U\simeq{}_{\ell}(A_U)$. Cela montre que l'on peut trouver un morphisme \inserted{\uline{fini}} \uline{étale} surjectif $S_1\to S$ qui splitte $ {}_{\ell}(A_U)$, \deleted{Remplaçant alors $S'$} et on peut supposer $S'$ \deleted{fini} \uline{étale} sur $S$, donc plat. Alors ($A'$ étant projectif sur $S'$) on \deleted{trouve une donnée de \unsure{des}} peut \uline{descendre} $A'$ en $A$, qui résoud le pb.

d) Cas $S$ réduit, \inserted{\deleted{la \unsure{normalisée}} tel qu'il existe $S'-S$ fini surjectif, avec $S'$ géom. unib.} $S'$ fini sur $S$. Utilisant le fait que $S$ est géom unibranche en les points de $S-U$, on voit que l'image inverse $U''=U'\times_U U'$ de $U$ dans $S''=S'\times_S S'$ est dense dans $S''$: \deleted{Par le cas \unsure{b|6}), on}

\diagram{\begin{tikzcd}[ampersand replacement=\&,column sep=small,row sep=small]
U \arrow[d,arrows={Hooks[right]-}] \& U' \arrow[l] \arrow[d,arrows={Hooks[right]-}] \& U'' \arrow[l,shift left=.45ex] \arrow[l,shift right=.45ex] \arrow[d,arrows={Hooks[right]-}]\\
S \& S' \arrow[l] \& S'' \arrow[l,dash,shift left=.45ex] \arrow[l,dash,shift right=.45ex]
\end{tikzcd}}

Par le cas c), on trouve sur $S'$ un schéma abélien $A'$, et par la pleine fidélité (pour $S'',U''$) une donnée de descente sur $A'$, relativement à $S'\to S$. Notons que le morphisme $S'-S$, étant un épimorphisme fini, est un morphisme de descente pour les schémas abéliens, \deleted{\unsure{\illegible{1}|E}} et aussi pour les gpes $p$-divisibles. Donc on est ramené à
% slanted
\margin{lemme à dégager.}
% archive sheet number: 12

\scanpage{14}
\mspage{-0.19-}
\[
\handwritten{\phi}(\handwritten{\alpha}X+\handwritten{\beta}Y,\ F_4(X,Y))=F_4(-\handwritten{\beta},\handwritten{\alpha}).
\]
Then it is easy to check that $(\handwritten{\mathbb A}^{1},\handwritten{\phi})$ is a geometric quotient of $X$ by $\mathrm{SL}(2)$. Finally, the morphism $\handwritten{\unsure{\Psi|\overline{\underline\psi}}}$ is not even closed. Let $Z$ be the closed subscheme of $G\times X$ whose closed points are the pairs:
\[
\begin{pmatrix}0&-\lambda^{-1}\\ \lambda&0\end{pmatrix}\times(\lambda X+Y,\ X^2Y^2)
\]
for $\lambda\in\handwritten{\mathbb C}-(0)$. The image under $\handwritten{\Psi}$ is the above point is:
\[
(-\lambda X+Y,\ X^2Y^2)\times(\lambda X+Y,\ X^2Y^2).
\]
In the closure of this set is the extra point:
\[
(Y,\ X^2Y^2)\times(Y,\ X^2Y^2).
\]

\uline{Lemma} \uline{0.5}: Let $\handwritten{\sigma}$ be an action of $\handwritten{\mathbb G}_m$ on an algebraic scheme $X$, all over a field $k$. Assume that $X$ admits an immersion in projective space $\handwritten{\mathbb P}_n$ for which the action $\handwritten{\sigma}$ extends to an action of $\handwritten{\mathbb G}_m$ on $\handwritten{\mathbb P}_n$. Then:
\[
\begin{aligned}
\handwritten{\sigma}\text{ is proper }\handwritten{\Longleftrightarrow}{}&
\left\{\begin{aligned}
&\text{(i) }\handwritten{\sigma}\text{ is separated,}\\
&\text{(ii) }S_1(X)=\handwritten{\phi}\,.
\end{aligned}\right.
\end{aligned}
\]

\uline{Proof}: The implication $\handwritten{\Longrightarrow}$ is clear. Conversely, assume $\handwritten{\sigma}$ is separated and $S_1(X)=\handwritten{\phi}$. To prove that the morphism $\handwritten{\Psi}$ is proper, we use the valuative criterion (EGA 2, \S7). Let $R$ be a valuation ring over $k$, and let $K$ be its quotient field. We may assume that $R$ contains an algebraically closed overfield of $k$, $\handwritten{\Omega}$, which is isomorphic to its residue field (Remark 7. 3. 9, EGA 2). Now suppose $\handwritten{\unsure{\rho|\phi}}\times\handwritten{\xi}$ is a $K$-valued but not $R$-valued point
% archive sheet number: 13
